\subsubsection{Test di Unità Frontend}

\renewcommand{\arraystretch}{1.2}
\begin{longtable}{|
		>{\raggedright\arraybackslash}p{0.15\columnwidth} |
		>{\raggedright\arraybackslash}p{0.325\columnwidth} |
		>{\raggedright\arraybackslash}p{0.325\columnwidth} |
		>{\centering\arraybackslash}p{0.10\columnwidth} |
	}
	\hline
	\rowcolor[gray]{0.9}
	\textbf{Codice Test$^G$} & \textbf{Descrizione} & \textbf{Valore atteso} & \textbf{Stato} \\
	\hline
	\endhead
	TU-B\_01 & Implementato da \textit{\seqsplit{test\_registration\_success}}; verifica che l'endpoint di registrazione dell'api possa essere chiamato e ritorni il valore una risposta positiva & \begin{itemize}[noitemsep, topsep=0pt, leftmargin=*]
		\item Risposta HTTP con stato 201
		\item Il campo "ok", della risposta in formato Json, è impostato a "True"
	\end{itemize} & S \\
	\hline
	\caption{Test di Unità Frontend}
\end{longtable}


ATTENZIONE QUELLO INSERITO È SOLO UN ESEMPIO PER LA TABELLA NIENTE DI PIÙ NIENTE DI MENO